\documentclass[a4paper,11pt]{article}
\usepackage[T1]{fontenc}
\usepackage{lmodern}
\usepackage[margin=22mm]{geometry}
\usepackage{amsmath,amssymb,mathtools}
\usepackage{titlesec,microtype}
\usepackage[hidelinks]{hyperref}
\setlength{\parindent}{0pt}
\setlength{\parskip}{4pt}
\setlength{\emergencystretch}{2em}
\titleformat{\section}{\normalfont\normalsize\bfseries}{}{0pt}{}
\titlespacing*{\section}{0pt}{10pt}{4pt}
\newcommand{\Prob}{\mathbb{P}}
\newcommand{\E}{\mathbb{E}}
\newcommand{\ind}{\mathbf{1}}
\newcommand{\R}{\mathbb{R}}
\DeclareMathOperator{\Dir}{Dirichlet}
\DeclareMathOperator*{\argmax}{arg\,max}
\hypersetup{pdftitle={Methods: point-configuration inference},pdfauthor={}}

\begin{document}
\textbf{Methods: point-configuration inference}

We jointly infer point configurations and annotator-specific count-error
probabilities from crowd annotations. A point configuration is a complete
collection of point coordinates, including their number; its points are not
selected independently. The method uses annotation coordinates and display
bounds, without reference labels or image intensities.

\section*{Response selection and candidate point configurations}

Annotations are mapped into parent-image coordinates and grouped into local
analysis regions. Pooled points are first partitioned spatially to prevent
long chains of nearby annotations from joining an entire image. Within each
partition, nearby points and points reported together are connected; their
components are enclosed in padded rectangles, and overlapping rectangles are
merged. These regions define units of inference, not biological object masks.
For each region and annotator, one available display is selected by greatest
region coverage, then by greatest edge margin. This avoids counting
overlapping displays as independent responses. Empty responses are retained;
unavailable responses contribute no evidence.

Let $Y_{rj}$ denote annotator $j$'s selected points in region $r$, $V_{rj}$ the
original display, and $J_r$ the available annotators; write
$Y_r=\{Y_{rj}:j\in J_r\}$. Each region has candidates
$\{S_{rk}\}_{k=1}^{K_r}$: the empty point configuration, distinct observed
responses, and one-center and two-center proposals from the pooled points.
The one-center proposal is their mean. The two-center proposal alternates
nearest-center assignment and mean updates, starting from the farthest pair.
Identical candidates are merged, but responses from different annotators are
retained. Point order is ignored and repeated clicks retain their multiplicity.
Candidates and regions remain fixed during estimation.

Write $T_r=k$ for the unknown candidate and $N_{rk}=|S_{rk}|$ for its count.
To avoid favoring counts represented by many similar candidates, the prior is
\begin{equation}
 a_{rk}=\Prob(T_r=k)=\frac{1}{|\mathcal N_r|K_{r,N_{rk}}},
 \qquad \mathcal N_r=\{N_{rk}\}_k,\qquad
 K_{rn}=\#\{k:N_{rk}=n\}.
 \label{eq:prior-config}
\end{equation}
Thus each represented count receives equal total prior mass. For partial
displays, the response model uses only the visible candidate points
$S_{rjk}^{\rm vis}=S_{rk}\cap V_{rj}$, preserving multiplicity. Define
$m_{rjk}=|S_{rjk}^{\rm vis}|$ and $n_{rj}=|Y_{rj}|$.

\section*{Model of annotation counts and positions}

Responses are conditionally independent given the region's point configuration.
Each annotator has a row-stochastic table $C_j\in(0,1)^{3\times3}$, with
$C_j[b,d]$ the probability of reporting count bin $d$ when the visible
underlying count is in bin $b$. The bins are $0$, $1$, and $2+$, represented by
$b(n)=\min(n,2)$. Each table has six free parameters and is shared across
regions. It describes count errors, not annotator-specific location precision.
The response likelihood is
\begin{equation}
 L_{rjk}(C_j)
 =C_j[b(m_{rjk}),b(n_{rj})]\,h_{m_{rjk}}(n_{rj})\,
 f_{m_{rjk},n_{rj}}(Y_{rj}\mid S_{rjk}^{\rm vis},V_{rj}).
 \label{eq:likelihood}
\end{equation}
Here $h_m$ allocates the $2+$ bin probability among exact counts:
\begin{equation}
 h_m(n)=
 \begin{cases}
  1,&n=0,1,\\
  q^{|n-u_m|}/Z_m,&n\geq2,
 \end{cases}
 \qquad u_m=\max(m,2),\quad
 Z_m=\frac{1+q-q^{u_m-1}}{1-q},\quad q=0.5.
 \label{eq:tail}
\end{equation}
Consequently, $C_j[b(m),b(n)]h_m(n)$ sums to one over all nonnegative counts;
counts above two are not discarded.

For a display $V$ with area $A_V$, define a Gaussian kernel truncated and
normalized on $V$, with a small uniform component:
\begin{equation}
 g_V(y;s)=(1-\varepsilon)
 \frac{\varphi_\sigma(y-s)}{\int_V\varphi_\sigma(u-s)\,du}
 +\frac{\varepsilon}{A_V},\qquad y\in V,
 \label{eq:kernel}
\end{equation}
and zero outside $V$. Here $\varphi_\sigma$ is the bivariate Gaussian density
with covariance $\sigma^2 I$, $\sigma=4$ pixels, and $\varepsilon=0.01$.
These values are fixed. Normalization uses the original display, not the
smaller analysis region.

For visible candidate points $S=\{s_1,\ldots,s_m\}$ and reported points
$Y=(y_1,\ldots,y_n)$, the spatial density $f_{m,n}$ distinguishes three
important response patterns. With $\bar s=m^{-1}\sum_a s_a$, they are
\begin{align}
 f_{m,1}(y\mid S,V)
 &=\frac{1}{2m}\sum_{a=1}^m g_V(y;s_a)
   +\frac12 g_V(y;\bar s), &&m>1,
 \label{eq:merge}\\
 f_{1,2}(y_1,y_2\mid\{s\},V)
 &=\frac12\bigl[g_V(y_1;s)g_V(y_2;2s-y_1)
 \notag\\[-2pt]
 &\hspace{2.8em}+g_V(y_2;s)g_V(y_1;2s-y_2)\bigr],
 \label{eq:split}\\
 f_{2,2}(y_1,y_2\mid S,V)
 &=\frac12\bigl[g_V(y_1;s_1)g_V(y_2;s_2)
 \notag\\[-2pt]
 &\hspace{2.8em}+g_V(y_2;s_1)g_V(y_1;s_2)\bigr].
 \label{eq:two-two}
\end{align}
Equation~\eqref{eq:merge} allows several points to be reported as one,
either near a retained point or near their centroid.
Equation~\eqref{eq:split} allows one point to be reported as two, with the
second near the reflection of the first through $s$.
Equation~\eqref{eq:two-two} places one response near each of two points.
The two response orders are averaged, so click ordering has no effect.
The mixture weights are fixed, not additional fitted parameters.

For an empty response, $f_{m,0}=1$; for no visible candidate points and
$n>0$, $f_{0,n}=A_V^{-n}$. All remaining cases with $m>0$ use
\begin{equation}
 f_{m,n}(Y\mid S,V)=\prod_{i=1}^{n}
       \left[\frac1m\sum_{a=1}^{m}g_V(y_i;s_a)\right].
 \label{eq:remaining}
\end{equation}
These are normalized, exchangeable coordinate densities conditional on count;
no additional factorial is included. Larger point configurations therefore
use an independent-draw approximation rather than all possible point matches.

\section*{Maximum-a-posteriori estimation}

Each row of $C_j$ has the fixed prior
$C_j[b,:]\sim\Dir(1+\eta_{b0},1+\eta_{b1},1+\eta_{b2})$, where
$\eta_{bd}=\lambda B_{bd}$, $\lambda=3$, and $B$ has diagonal entries $0.8$
and off-diagonal entries $0.1$. The diagonal preference is a modeling
assumption, not a supplied correct label. Up to a constant, the objective is
\begin{equation}
 \mathcal J(C)=\sum_r\log\!\left[
       \sum_k a_{rk}\prod_{j\in J_r}L_{rjk}(C_j)\right]
       +\sum_{j,b,d}\eta_{bd}\log C_j[b,d].
 \label{eq:objective}
\end{equation}
This is the response log likelihood, marginalized over candidate point
configurations, plus the log prior. Regions and candidates were constructed
from the data and are conditioned on; their construction is not jointly
modeled as part of a whole-image likelihood.

We maximize Eq.~\eqref{eq:objective} by expectation--maximization.
Given $C^{(t)}$, the E-step computes
\begin{equation}
 w_{rk}^{(t)}=\Prob(T_r=k\mid Y_r,C^{(t)})
 =\frac{a_{rk}\prod_{j\in J_r}L_{rjk}(C_j^{(t)})}
 {\sum_\ell a_{r\ell}\prod_{j\in J_r}L_{rj\ell}(C_j^{(t)})}.
 \label{eq:posterior}
\end{equation}
These weights provide the expected count-bin pairs for each annotator,
followed by an exact row-wise M-step:
\begin{align}
 A_{jbd}^{(t)}&=\sum_{r:j\in J_r}\sum_k w_{rk}^{(t)}
   \ind\{b(m_{rjk})=b\}\ind\{b(n_{rj})=d\},
 \label{eq:expected}\\
 C_j^{(t+1)}[b,d]&=
 \frac{A_{jbd}^{(t)}+\eta_{bd}}
 {\sum_{d'=0}^2 A_{jbd'}^{(t)}+\lambda}.
 \label{eq:update}
\end{align}
Thus annotator tables alter the evidence for candidate counts, and candidate
probabilities update the annotator tables. The added terms $\eta_{bd}$ are MAP
pseudocounts, not Dirichlet shape parameters. An unsupported row returns to
$B_{b,:}$. Candidate probabilities are derived quantities, not separate fitted
parameters; no additional annotator confidence score is applied to the output.

Calculations use log densities and log-sum-exp. We fit from three starts:
$C_j=B$, uniform rows, and rows sampled from the prior. Iteration stops when
the maximum absolute change in both $C$ and $w$ is below $10^{-7}$ and
$|\Delta\mathcal J|<10^{-7}\max(1,|\mathcal J|)$, with at most 2000 iterations.
The highest-objective converged fit is retained, and its posterior is
recomputed using the final tables. No output is accepted from an unconverged
fit. Multiple starts reduce, but do not eliminate, sensitivity to local optima.

\section*{Application and output}

For new images, annotator tables are held fixed and only the candidate
posterior in Eq.~\eqref{eq:posterior} is evaluated. An annotator absent from
training receives the arithmetic mean of the fitted tables. The exact count
distribution, expected count, and selected point configuration are
\begin{align}
 \Prob(N_r=n\mid Y)&=\sum_{k:N_{rk}=n}w_{rk},
 &\E[N_r\mid Y]&=\sum_k w_{rk}N_{rk},
 \label{eq:count-output}\\
 \widehat S_r&=S_{r,k_r^*},
 &k_r^*&=\argmax_k w_{rk}.
 \label{eq:map-output}
\end{align}
The selected candidate retains its complete point combination. A spatial
summary is the expected point-mass measure
\begin{equation}
 M_r(dx)=\sum_k w_{rk}\sum_{s\in S_{rk}}\delta_s(dx),
 \qquad M_r(\R^2)=\E[N_r\mid Y],
 \label{eq:mass}
\end{equation}
where $\delta_s$ is unit mass at $s$. For visualization, bilinear deposition
and a normalized Gaussian filter with standard deviation 1.5 pixels preserve total mass;
smoothing does not alter posterior weights. Peaks from alternative point
configurations may coexist in this mean map without forming one jointly
supported point configuration.

Inference is conditional on the candidate set, fixed spatial model, and fitted
annotator tables. It does not integrate uncertainty in preprocessing or table
estimation. Entirely unannotated regions are not proposed, and point
configurations outside the candidate set cannot be recovered. The method
estimates count-error tendencies without reference labels, but does not learn
annotator-specific spatial accuracy. Its statistical estimates depend on the
stated assumptions and are not, by themselves, evidence of biological truth.

\end{document}
